\chapter{Algorytmy}

W tym rozdziale przedstawiamy dwa kluczowe algorytmy dla markowskich procesów decyzyjnych z quasi-hiperbolicznym dyskontowaniem: algorytm oceny polityki oraz algorytm QH Q-Learning. Oba algorytmy są bezmodelowe (model-free), co oznacza, że nie wymagają znajomości dynamiki środowiska $q(s' \mid s, a)$.

\section{Algorytm 1: Bezmodelowa ocena polityki dla dyskontowania QH}

Pierwszy algorytm pozwala na ocenę wartości zadanej polityki bez znajomości modelu środowiska.

\subsection{Wyprowadzenie wzoru dla oceny polityki}

Z twierdzenia o redukcji polityk (Twierdzenie~\ref{thm:policy-reduction}):
\begin{equation}
J^{\alpha, \beta}_{\mu, \phi_s}(s) = \sum_a \mu(a\mid s) \left[ r(s, a) + \alpha\beta \sum_{s'} q(s'\mid s, a) \, J^{\beta}_{\phi_s}(s') \right].
\label{eq:policy-eval-start}
\end{equation}

Celem jest wyrażenie zależności wyłącznie przez $J^{\alpha, \beta}_{\phi_s}(s')$, bez jawnego $J^{\beta}$.

\subsubsection{Definicje dla polityki stacjonarnej}

Dyskontowanie wykładnicze:
\begin{equation}
J^{\beta}_{\phi_s}(s') = \mathbb{E}_{\phi_s} [r(s', \cdot)] + \beta \mathbb{E}_{\phi_s, q} [J^{\beta}_{\phi_s}(s'') \mid s'],
\label{eq:exp-value}
\end{equation}
gdzie $\mathbb{E}_{\phi_s} [r(s', \cdot)] = \sum_{a'} \phi_s(a'\mid s') \, r(s', a')$.

Quasi-hiperboliczne:
\begin{equation}
J_{\phi_s}(s') = \mathbb{E}_{\phi_s} [r(s', \cdot)] + \alpha \beta \mathbb{E}_{\phi_s, q} [J^{\beta}_{\phi_s}(s'') \mid s'].
\label{eq:qh-value-stat}
\end{equation}

\subsubsection{Izolacja $\alpha \beta J^{\beta}_{\phi_s}(s')$}

Pomnóżmy równanie~\eqref{eq:exp-value} przez $\alpha \beta$:
\begin{equation}
\alpha \beta J^{\beta}_{\phi_s}(s') = \alpha \beta \mathbb{E}_{\phi_s} [r(s', \cdot)] + \alpha \beta^2 \mathbb{E}_{\phi_s, q} [J^{\beta}_{\phi_s}(s'') \mid s'].
\label{eq:mult-ab}
\end{equation}

Pomnóżmy równanie~\eqref{eq:qh-value-stat} przez $\beta$:
\begin{equation}
\beta J_{\phi_s}(s') = \beta \mathbb{E}_{\phi_s} [r(s', \cdot)] + \alpha \beta^2 \mathbb{E}_{\phi_s, q} [J^{\beta}_{\phi_s}(s'') \mid s'].
\label{eq:mult-b}
\end{equation}

Odejmując równanie~\eqref{eq:mult-b} od równania~\eqref{eq:mult-ab}:
\begin{equation}
\alpha \beta J^{\beta}_{\phi_s}(s') - \beta J_{\phi_s}(s') = \alpha \beta \mathbb{E}_{\phi_s} [r(s', \cdot)] - \beta \mathbb{E}_{\phi_s} [r(s', \cdot)].
\end{equation}

Przekształcamy:
\begin{equation}
\alpha \beta J^{\beta}_{\phi_s}(s') = \beta J_{\phi_s}(s') - (1 - \alpha) \beta \mathbb{E}_{\phi_s} [r(s', \cdot)].
\label{eq:isolated}
\end{equation}

Podstawiając równanie~\eqref{eq:isolated} do równania~\eqref{eq:policy-eval-start}:
\begin{align}
J^{\alpha, \beta}_{\mu, \phi_s}(s) &= \sum_a \mu(a\mid s) \left[ r(s, a) + \sum_{s'} q(s'\mid s, a) \left[ -(1 - \alpha)\beta \sum_{a'} \phi_s(a'\mid s') r(s', a') + \beta J_{\phi_s}(s') \right] \right].
\label{eq:final-policy-eval}
\end{align}

\subsection{Opis algorytmu}

\begin{algorithm}[H]
\caption{Bezmodelowa ocena polityki dla dyskontowania QH}
\label{alg:policy-eval}
\begin{algorithmic}[1]
\REQUIRE Kroki uczenia $(\eta_n)$, $(\theta_n)$; współczynniki dyskontowania $\alpha$, $\beta$; polityka próbkowania $\nu$; polityki oceniane $\mu$, $\phi_s$
\STATE \textbf{Inicjalizacja:} $J_0, W_0 \in \mathbb{R}^{|S|}$
\FOR{$n = 0, 1, 2, \dots$}
\FOR{$s \in S$}
\STATE Próbkuj $a \sim \nu(\cdot|s)$
\STATE Zaobserwuj nagrodę $r(s, a)$ oraz stan kolejny $s' \sim q(\cdot|s, a)$
\STATE Próbkuj $a' \sim \phi_s(\cdot|s')$
\STATE Zaobserwuj nagrodę $r(s', a')$
\STATE Oblicz cel TD: $r^{\text{target}}_n(s) \leftarrow r(s, a) - (1 - \alpha)\beta \, r(s', a') + \beta W_n(s')$
\STATE Zaktualizuj $W$: $W_{n+1}(s) \leftarrow W_n(s) + \eta_n \left[ \frac{\phi_s(a|s)}{\nu(a|s)} r^{\text{target}}_n(s) - W_n(s) \right]$
\STATE Zaktualizuj $J$: $J_{n+1}(s) \leftarrow J_n(s) + \theta_n \left[ \frac{\mu(a|s)}{\nu(a|s)} r^{\text{target}}_n(s) - J_n(s) \right]$
\ENDFOR
\ENDFOR
\STATE \textbf{Wynik:} $J_n$ zbieżne do $J^{\alpha, \beta}_{\mu, \phi_s}$; $W_n$ zbieżne do $J^{\alpha, \beta}_{\phi_s}$
\end{algorithmic}
\end{algorithm}

\subsection{Konwergencja algorytmu}

\begin{theorem}[Konwergencja Algorytmu~\ref{alg:policy-eval}]
\label{thm:policy-eval-conv}
Jeżeli kroki uczenia $\eta_n$, $\theta_n$ spełniają warunki Robbins--Monro:
\[
\sum_n \eta_n = \infty, \quad \sum_n \eta_n^2 < \infty \quad \text{oraz} \quad \sum_n \theta_n = \infty, \quad \sum_n \theta_n^2 < \infty,
\]
oraz dodatkowo $\theta_n / \eta_n \to 0$ (dwie skale czasowe), to ciągi $(W_n, J_n)$ są zbieżne prawie na pewno do $(J^{\alpha, \beta}_{\phi_s}, J^{\alpha, \beta}_{\mu, \phi_s})$.
\end{theorem}

Dowód opiera się na teorii aproksymacji stochastycznej o dwóch skalach czasowych i jest analogiczny do dowodu Twierdzenia~\ref{thm:qh-qlearning-conv}.

\section{Algorytm 2: QH Q-Learning (Bezmodelowy)}

Drugi algorytm pozwala na znalezienie optymalnej polityki bez znajomości modelu środowiska.

\subsection{Relacja funkcji Q}

Kluczowa relacja między funkcjami Q dla różnych dyskontowań:
\begin{equation}
Q^{\alpha, \beta}_{\mu, \phi_s}(s,a) = (1 - \alpha) r(s,a) + \alpha Q^{\beta}_{\phi_s}(s,a), \qquad Q^{\alpha, \beta}_*(s,a) = (1 - \alpha) r(s,a) + \alpha Q^{\beta}_*(s,a).
\label{eq:q-relation}
\end{equation}

\subsection{Opis algorytmu}

\begin{algorithm}[H]
\caption{QH Q-Learning (Bezmodelowy)}
\label{alg:qh-qlearning}
\begin{algorithmic}[1]
\REQUIRE Kroki uczenia $(\eta_n)$, $(\theta_n)$; współczynniki dyskontowania $\alpha$, $\beta$
\STATE \textbf{Inicjalizacja:} $W_0, Q_0 \in \mathbb{R}^{|S| \times |A|}$
\FOR{$n = 0, 1, 2, \ldots$}
\FOR{$(s, a) \in S \times A$}
\STATE Próbkuj $s' \sim q(\cdot\mid s, a)$ i zaobserwuj $r(s, a)$
\STATE $W_{n+1}(s, a) \leftarrow W_n(s, a) + \eta_n \left[ r(s, a) + \beta \max_{a'} W_n(s', a') - W_n(s, a) \right]$
\STATE $Q_{n+1}(s, a) \leftarrow Q_n(s, a) + \theta_n \left[ (1 - \alpha) r(s, a) + \alpha W_n(s, a) - Q_n(s, a) \right]$
\ENDFOR
\ENDFOR
\STATE \textbf{Wynik:} $Q_n$ (QH) oraz $W_n$ (wykładnicze); $\mu^\star(s) = \arg\max_a Q_n(s, a)$, $\phi_s^\star(s) = \arg\max_a W_n(s, a)$
\end{algorithmic}
\end{algorithm}

\subsection{Konwergencja algorytmu}

Zbieżność algorytmu jest zagwarantowana przez Twierdzenie~\ref{thm:qh-qlearning-conv}:

\begin{itemize}
\item \textbf{Szybka skala ($W_n$):} Klasyczny Q-learning prowadzi do $Q^{\beta}_{\phi_s^\star}$.
\item \textbf{Wolna skala ($Q_n$):} Układ liniowy ze stanem quasi-stałym $W_n \approx Q^{\beta}_{\phi_s^\star}$ zbiega do $Q^{\alpha, \beta}_{\mu^\star, \phi_s^\star}$.
\item Zachowanie greedy zapewnia odzyskanie optymalnych polityk.
\end{itemize}

\subsection{Złożoność obliczeniowa}

Analizujemy złożoność czasową i pamięciową algorytmów:
\begin{itemize}
\item \textbf{Pamięć:} $O(|S| \times |A|)$ dla tablic $W$ i $Q$
\item \textbf{Czas na iterację:} $O(|S| \times |A|)$ dla pełnej aktualizacji
\item Algorytmy dwuskalowe wymagają większej liczby iteracji dla zbieżności, ale oferują gwarancje konwergencji
\end{itemize}

\section{Porównanie z metodami tradycyjnymi}

Algorytm QH Q-Learning różni się od klasycznego Q-learningu w następujących aspektach:

\begin{itemize}
\item \textbf{Dwie funkcje Q:} QH Q-Learning utrzymuje dwie funkcje Q: $W$ dla dyskontowania wykładniczego i $Q$ dla dyskontowania quasi-hiperbolicznego.
\item \textbf{Dwie skale czasowe:} Aktualizacje $W$ i $Q$ odbywają się z różnymi krokami uczenia, gdzie $\theta_n / \eta_n \to 0$.
\item \textbf{Polityka dwuskładnikowa:} Wynikiem jest para $(\mu^\star, \phi_s^\star)$ zamiast pojedynczej polityki.
\item \textbf{Niespójność czasowa:} Optymalna polityka nie jest spójna w czasie, co wymaga specjalnej interpretacji i zastosowania.
\end{itemize}
